\subsection{万有纤维丛}

设 G 是一个拓扑群，B 和 B' 是拓扑空间，且 \((\mathrm{E}, p)\) 和 \((\mathrm{E}', p')\) 分别是结构群为 G、底空间分别为 B 和 B' 的主纤维丛。

设 U 是 B 的一个开子集，且 \(f'\colon \overline{p}^{-1}(U) \to E'\) 是一个连续映射，它分别与 \(\overline{p}^{-1}(U)\) 和 \(E'\) 中 G 的运算相容。因而存在唯一连续映射 \(f\colon U \to B'\)，使得 \(f \circ p_{U} = p' \circ f'\)，并且交换方形

\[
\begin{array}{c} \mathrm {E_ {U}} \xrightarrow {f ^ {\prime}} \mathrm {E^ {\prime}} \\ \Biggl \downarrow p _ {\mathrm{U}} \\ \mathrm{U} \xrightarrow {f} \mathrm {B^ {\prime}} \end{array}
\]

随后该交换方形是笛卡尔的（I, p. 94, example (FP)）。

对于 B 的每个开子集 U，记 \(\mathcal{F}(\mathrm{U})\) 为所有连续映射 \(g\colon \mathrm{E}_{\mathrm{U}}\to \mathrm{E}'\) 的集合；这些映射分别与 \(\overline{p}^{-1}(\mathrm{U})\) 和 \(\mathrm{E}'\) 中 G 的运算相容。对于 B 的每一对开集 (U,V)，若 \(\mathrm{U}\subset \mathrm{V}\)，令 \(r_{\mathrm{UV}}\colon \mathcal{F}(\mathrm{V})\to \mathcal{F}(\mathrm{U})\) 为由 \(r_{\mathrm{UV}}(g) = g\mid \mathrm{E}_{\mathrm{U}}\) 定义的映射。立即可验证，这样便在 B 上定义了层 \(\mathcal{F} = ((\mathcal{F}(\mathrm{U})),(r_{\mathrm{UV}}))\)。我们称此层为 B 上由 E 到 \(\mathrm{E}'\) 的群 G 主纤维丛态射层。

命题 4. --- 若空间 B 是仿紧的，且空间 E' 与一个点同伦等价，则 B 上由 E 到 E' 的群 G 主纤维丛态射层是柔软层。

存在 B 的一个开覆盖 \((\mathrm{U}_{j})_{j\in\mathrm{J}}\)，使得对于每个 \(j\in J\)，纤维丛 \(E_{U_{j}}\) 可平凡化。由于空间 B 是仿紧的，可以设覆盖 \((\mathrm{U}_{j})_{j\in\mathrm{J}}\) 是局部有限的（TG, IX, p. 49），并选取 B 的一个覆盖 \((\mathrm{A}_{j})_{j\in\mathrm{J}}\)，使得对于每个 \(j\in J\)，集合 \(A_{j}\) 在 B 中是闭的且包含于 \(U_{j}\) 中（TG, IX, p. 49, prop. 4 and p. 48, cor. 1）。

由 I, p. 65, cor. 2 of prop. 6，要证明命题只需建立以下断言：设 U 是 B 的一个开子集，使得主纤维丛 \((E_{U}, p_{U})\) 可平凡化；设 A 是包含于 U 中的 B 的闭子集；设 V 是包含 A 且包含于 U 中的开邻域，并令 f 为 \(\mathcal{F}(V)\) 中的一个元素；则存在 V 中 A 的一个开邻域 W，以及一个元素 \(f'\)，它属于 \(\mathcal{F}(U)\)，满足 \(r_{\mathrm{WU}}(f') = r_{\mathrm{WV}}(f)\)。我们来证明这个断言。

取 B 的一个开子集 W，使得 \(A \subset W \subset \overline{W} \subset V\)。令 \(s: U \to E_{U}\) 为 \((\mathrm{E}_{\mathrm{U}}, p_{\mathrm{U}})\) 的一个截面。将 Corollary 3 (IV, p. 438) 应用于空间 B、闭集 \(\overline{W}\)、\(\overline{W}\) 的邻域 V，以及映射 \(g = f \circ (s\textbar{} V)\)，该映射从 \(V\) 到 \(E'\)。因而存在连续映射 \(\widetilde{g}: B \to E'\)，它在 \(\overline{W}\) 上与 g 重合。令 \(h = \widetilde{g}\textbar{} U\)。有 \(h\textbar{} W = g\textbar{} W = f \circ (s\textbar{} W)\)。

映射 H: U × G → E' 由 H(x, g) = h(x) · g 定义，对于 (x, g) ∈ U × G，它是连续的，并且分别与 U × G 和 E' 中 G 的运算相容。置 \(f' = H \circ s^{-1}\)；它是 \(\mathcal{F}(U)\) 中的一个元素。对于所有 \(x \in W\)，映射 f 和 \(f'\) 在点 \(s(x)\) 处重合，因此由于它们是主纤维丛态射，在 \(\overline{p}^{-1}(x)\) 的每一点处也重合。命题得证。

定理 1. --- 设 G 是一个拓扑群，设 \(B_{u}\) 是一个拓扑空间，且 \((\mathrm{E}_{\mathrm{u}}, p_{\mathrm{u}})\) 是以 \(B_{u}\) 为底、群 G 的主纤维丛。我们假设空间 \(E_{u}\) 与一个点同伦等价。

设 B 是一个仿紧拓扑空间。

\begin{enumerate}
\def\labelenumi{\alph{enumi})}
\item
每个群 G、底空间为 B 的主纤维丛，都同构于某个形如 \(f^{*}E_{u}\) 的主纤维丛，其中 \(f: B \to B_{u}\) 是连续映射。
\item
令 \(f_{0}\) 和 \(f_{1}\) 为从 B 到 \(B_{u}\) 的连续映射。要使 \(f_{0}^{*}E_{u}\) 和 \(f_{1}^{*}E_{u}\) 成为以 B 为底、群 G 的同构主纤维丛，\(f_{0}\) 和 \(f_{1}\) 同伦是充要条件。
\end{enumerate}

换言之，存在从 \([B, B_{u}]\) 到 \(P(B; G)\) 的映射，它将从 B 到 \(B_{u}\) 的连续映射 f 的同伦类对应到主纤维丛 \(f^{*}E_{u}\) 的同构类。该映射是双射。

设 E 是群 G、底空间为 B 的主纤维丛。由 Prop. 4，B 上的主纤维丛态射层 \(\mathcal{F}\)（由 E 到 \(\mathrm{E_u}\)）是柔软的。因此，\(\mathcal{F}(\mathrm{B})\) 非空，从而 \(a\)。

令 \(f_{0}\) 和 \(f_{1}\) 为从 B 到 \(B_{u}\) 的连续映射。若映射 \(f_{0}\) 和 \(f_{1}\) 同伦，则主纤维丛 \(f_{0}^{*}E_{u}\) 和 \(f_{1}^{*}E_{u}\) 同构（IV, p. 445, cor. 2）。现在证明逆命题。对于 \(\alpha\in\{0,1\}\)，记 \(E_{\alpha}\) 为以 \(B_{u}\) 为底的主纤维丛 \(f_{\alpha}^{*}E_{u}\)，记 \(g_{\alpha} \colon E_{\alpha} \to E_{u}\) 为第一投影，并记 \(p_{\alpha} \colon E_{\alpha} \to B\) 为第二投影。设 \(i \colon E_{0} \to E_{1}\) 是主纤维丛同构。令 \(p\) 为映射 \(p_{0} \times Id_{I} \colon E_{0} \times I \to B \times I\)。

由于空间 B × I 是仿紧的（TG, IX, p. 70, prop. 17），B × I 上从 E₀ × I 到 Eᵤ 的主纤维丛态射层 G 是柔软的（IV, p. 446, prop. 4）。置 A = B × \{0,1\}，U = B × ([0, ½[∪]½, 1])，并通过下式定义 G(U) 中的元素 g：

\[
g (x, t) = \left\{ \begin{array}{l l} g _ {0} (x) & \text { for } (x, t) \in \mathrm{E} _ {0} \times [ 0, \frac {1}{2} [, \\ g _ {1} \circ i (x) & \text { for } (x, t) \in \mathrm{E} _ {0} \times ] \frac {1}{2}, 1 ]. \end{array} \right.
\]

由于层 \(\mathcal{G}\) 是柔软的，存在一个 \(h \in \mathcal{G}(\mathrm{B} \times \mathbf{I})\) 以及 A 在 U 中的一个开邻域 V，使得 \(h \mid \mathrm{V} = g \mid \mathrm{V}\)；这样的元素是一个连续映射 \(\mathrm{H} \colon \mathrm{E}_0 \times \mathbf{I} \to \mathrm{E}_{u}\)，它与 G 的运算相容，并满足 \(\mathrm{H}(x,0) = g_0(x)\)、\(\mathrm{H}(x,1) = g_1(i(x))\)，对于所有 \(x \in \mathrm{E}_0\) 都成立。因而存在映射 \(h' \colon \mathrm{B} \times \mathbf{I} \to \mathrm{B}_{u}\)，使得 \(h'(p_0(x),t) = p_{\mathrm{u}}(\mathrm{H}(x,t))\)，对于 \(x \in \mathrm{E}_0\) 和 \(t \in \mathbf{I}\) 成立；该映射连续，并且是连接 \(f_0\) 与 \(f_1\) 的同伦。

推论。 --- 设 G 是一个拓扑群，B 和 B' 是仿紧拓扑空间。设 (E,p) 和 (E',p') 分别是群 G、底空间为 B 和 B' 的主纤维丛。假设空间 E 和 E' 都与一个点同伦等价。则空间 B 和 B' 同伦等价。

事实上，存在连续映射 \(f \colon \mathrm{B} \to \mathrm{B}'\)，使得主纤维丛 E 同构于主纤维丛 \(f^* \mathrm{E}'\)，以及连续映射 \(g \colon \mathrm{B}' \to \mathrm{B}\)，使得主纤维丛 \(\mathrm{E}'\) 同构于主纤维丛 \(g^* \mathrm{E}\)（定理 1, (a)）。于是主纤维丛 \((g \circ f)^* \mathrm{E}\) 与 E 同构，因此映射 \(g \circ f\) 同伦于映射 \(\mathrm{Id}_{\mathrm{B}}\)（定理 1, (b)）。同样，映射 \(f \circ g\) 同伦于映射 \(\mathrm{Id}_{\mathrm{B}'}\)。

设 G 是一个拓扑群，B 是一个拓扑空间，E 是以 B 为底的主 G-丛。若空间 E 与一个点同伦等价，则称主丛 E 是具有仿紧底空间的主 G-丛的万有丛，并称空间 B 是 G 的分类空间。若 G 有两个分类空间是仿紧的，则它们同伦等价。当分类空间存在时，可以把具有仿紧底空间的主 G-丛的同构类的研究看作同伦论中的一个问题。

例。 --- 赋予 R 映射 \(p: R \to S^{1}\)、\(t \mapsto e^{2\pi it}\)，并赋予 Z 平移作用后，空间 R 是群 Z 的主覆盖。因此，空间 \(S^{1}\) 是群 Z 的分类空间。

对于每个离散群 G，我们将在下一节构造一个度量空间，它是 G 的分类空间。

